% Generated draft: requires mathematical and source review.
\documentclass[11pt]{article}
\usepackage[margin=25mm]{geometry}
\usepackage{fontspec}
\setmainfont{Latin Modern Roman}
\usepackage{amsmath,amssymb,mathtools,unicode-math}
\setmathfont{Latin Modern Math}
\usepackage{enumitem,needspace,xurl,hyperref}
\setlength{\parindent}{0pt}
\setlength{\parskip}{5pt}
\setlength{\emergencystretch}{3em}
\newcommand{\N}{\mathbb{N}}
\newcommand{\Z}{\mathbb{Z}}
\newcommand{\Q}{\mathbb{Q}}
\newcommand{\R}{\mathbb{R}}
\newcommand{\C}{\mathbb{C}}
\newcommand{\F}{\mathbb{F}}
\newcommand{\A}{\mathbb{A}}
\newcommand{\PP}{\mathbb{P}}
\newcommand{\E}{\mathbb{E}}
\newcommand{\eps}{\varepsilon}
\newcommand{\dd}{\,\mathrm{d}}
\begin{document}
\section*{UnsolvedMath ID: 30005710}
\textbf{Source problem code:} catalogue entry
\subsection{Definitions and mathematical statement}
Let \(d,j,n,k\in\N\) be positive integers, and let \(\mathcal C=\{C_1,\ldots,C_j\}\) be the supplied collection of proper convex bodies in \(\R^d\), with \(0\in\operatorname{int}_{\R^d}(C_r)\) for every \(r\). Under the usual convention, a convex body is a compact convex set with nonempty interior; the record does not explain whether the adjective ``proper'' imposes any additional condition. Let \(\mathcal P\) be a prescribed class of convex partitions of \(\R^k\).

For a \(k\)-dimensional linear subspace \(L\subseteq\R^d\), write \(\operatorname{vol}_L\) for the \(k\)-dimensional Lebesgue measure induced by its Euclidean structure. To make the volume requirement explicit, a partition up to null boundaries consists of measurable convex sets \(A_1,\ldots,A_n\subseteq L\) satisfying
\[
L=\bigcup_{i=1}^n A_i,
\qquad
\operatorname{vol}_L(A_i\cap A_h)=0\quad(i\ne h).
\]
The supplied record does not specify whether its partitions use this convention, literal disjointness, closed cells with common boundaries, or additional cell conditions.

A further convention is needed to interpret membership in \(\mathcal P\) for partitions of \(L\), since \(\mathcal P\) is defined on \(\R^k\). One possible precise interpretation is that there is a linear isometry \(T\colon\R^k\to L\) such that
\[
\bigl(T^{-1}(A_1),\ldots,T^{-1}(A_n)\bigr)\in\mathcal P.
\]
This interpretation is conditional: the record does not state whether the identification is existential, prescribed, or immaterial because \(\mathcal P\) is invariant under orthogonal changes of coordinates.

\textbf{Equipartition question, subject to these unresolved conventions.} Given the above data, do there exist a \(k\)-dimensional linear subspace \(L\subseteq\R^d\) and an admissible convex partition \((A_1,\ldots,A_n)\) of \(L\), belonging to the prescribed class after the intended identification, such that, simultaneously for every \(r\in\{1,\ldots,j\}\),
\[
\operatorname{vol}_L(C_r\cap A_1)=\cdots=
\operatorname{vol}_L(C_r\cap A_n)
=\frac{1}{n}\operatorname{vol}_L(C_r\cap L)?
\]
Here \(A_i\subseteq L\), so \(C_r\cap A_i=C_r\cap L\cap A_i\). The final equality follows from the partition convention above and finite volume; it makes explicit the equal-volume requirement in the supplied statement.

A necessary dimensional condition is \(k\le d\). The supplied range only says \(d,j,n,k\ge1\); if \(k>d\), the requested subspace does not exist. Moreover, an arbitrary prescribed class \(\mathcal P\) need not admit a solution. The source supplies no restrictions on \(\mathcal P\) or numerical hypotheses that would turn this question into an unconditional existence assertion. Such restrictions must therefore remain unresolved rather than being inferred from the title or surrounding discussion of Euclidean vector bundles.
\subsection{Short English statement}
The question asks whether one can choose a linear \(k\)-dimensional section through the origin and divide that section into \(n\) convex cells of a prescribed type so that all \(j\) convex bodies have equal sectional volume in every cell. Precision requires clarification of the admissible partition class, its transport to each subspace, and boundary conventions. Existence is impossible for \(k>d\), and the record gives no general hypotheses guaranteeing it for arbitrary \(\mathcal P\). The supplied partially solved status is dated curation, not an independently verified conclusion.
\subsection{Sources}
\begin{itemize}
\item Supplied local UnsolvedMath registry record: permanent ID \(30005710\), title ``Convex Equipartitions in Tautological Vector Bundles,'' source problem code \texttt{catalogue entry}, collection ``Oberwolfach Reports Open Problems.'' The record gives proposed year 2023 and no named proposer.
\item Supplied primary-source citation: \emph{Geometric, Algebraic and Topological Combinatorics} (2024), \emph{Oberwolfach Reports}, pp.~3249--3318, DOI \texttt{10.4171/owr/2023/58}. Supplied source URL and literature-source link labelled ``original Oberwolfach report'': \url{https://doi.org/10.4171/owr/2023/58}. The differing year indications are retained as supplied.
\item Supplied literature-source attribution: Blagojevic--Crabb, \emph{Many partitions of mass assignments}, \url{https://doi.org/10.4171/DM/980}. No further bibliographic details are supplied.
\item Literature assessment dated \texttt{2026-08-22}: supplied dated curation, not independently verified. The record labels the problem \texttt{partially\_solved} and reports positive theorems for important orthant and generalized-Voronoi settings, while describing its unrestricted formulation as unresolved. Its background explicitly calls this a dated literature triage derived from supplied Oberwolfach curation and requests independent verification of time-sensitive status and cited claims.
\item The registry labels the statement \texttt{corrected\_verified} and says that it was reconstructed and checked against the official Oberwolfach report. This is a supplied verification claim; no independent source visit or verification is asserted here.
\item The damaged original excerpt attributes surrounding partition constructions to Gromov; Palic with Blagojevic and Karasev, internal reference [15]; Blagojevic and Soberon, [9]; and Blagojevic and Sadovek, [8]. It discusses generalized Voronoi sites, configuration spaces, iterated partitions, sphere-waist results, and Nandakumar--Ramana-Rao problems. Complete bibliographic entries and the hypotheses needed to apply those constructions are absent from the supplied record.
\end{itemize}
\end{document}
